\documentclass[11pt]{article}
\usepackage{amsmath,amssymb,amsthm,hyperref}
\newtheorem{theorem}{Theorem}
\newtheorem{lemma}{Lemma}
\newtheorem{corollary}{Corollary}
\newcommand{\E}{\mathbb E}
\title{A universal endpoint profile for all comparison columns}
\author{Asymptotic research note; independently reviewed}
\date{30 September 2026}
\begin{document}\maketitle
For every $m$, let $p_1,\ldots,p_m:[K]\to[0,1]$ be nondecreasing,
with positive probability weights $w_q$, and suppose
\[
 \E_w\prod_{j\in S}p_j=\frac1{|S|+1}\qquad(S\subseteq[m]).
\]
The number of rows and weights can depend on $m$.
Put $A=\prod_jp_j$, $\eta(q)=(m+1)w_qA(q)$ and, on $A>0$,
\[
 r_j=\frac{1-p_j}{p_j},\qquad \lambda=\sum_jr_j.
\]
The separately proved endpoint theorem in
\path{../../independent_directions/endpoint_poisson_rigidity.tex}
shows that $\lambda$ under $\eta$ converges to $\operatorname{Exp}(1)$,
and that $\eta(\max_jr_j\ge\varepsilon)\le1/(m\varepsilon)$.
The present note strengthens that conclusion from the sum of the odds
to every individual column, uniformly.

\begin{lemma}[A fixed number of endpoint failures]\label{gamma}
Fix an integer $\ell\ge0$, and choose any $\ell$ columns $J=J_m$.
Define a probability measure on the original rows by
\[
 \nu_J(q)=(m+1)\binom m\ell w_q
       \prod_{j\in J}(1-p_j(q))\prod_{i\notin J}p_i(q).
\]
On its support set $\lambda_{-J}=\sum_{i\notin J}(1-p_i)/p_i$.
As $m\to\infty$, its distribution under $\nu_J$ converges to the
Gamma law of shape $\ell+1$ and rate one. The convergence is uniform
in the choice of $J$, in the sequential sense that it holds for every
sequence $J_m$.
\end{lemma}
\begin{proof}
A particular failure pattern of size $s$ has probability
$1/((m+1)\binom ms)$, by expanding the mixed moments.
Writing $d=m-\ell$, the exact generating polynomial for additional
failures is consequently
\[
 (m+1)\binom m\ell\E_w\left[
  \prod_{j\in J}(1-p_j)
  \prod_{i\notin J}(p_i+z(1-p_i))\right]
 =\sum_{s=0}^{d}\binom{\ell+s}{\ell}z^s.
\]
For any one index $i\notin J$, replace its factor on the left by
$1-p_i$. The resulting coefficient of $z^s$ is
\[
 \frac{s+1}{d}\binom{\ell+s+1}{\ell}.
\]
Thus, for $0\le z<1$, this polynomial is at most
$(\ell+1)/(d(1-z)^{\ell+2})$.

The event that an odds ratio outside $J$ is at least $\varepsilon$
is an initial segment. Choosing a column witnessing this event at
its last row bounds its contribution to the first generating
polynomial by
\[
 (\varepsilon^{-1}+z)
 \frac{\ell+1}{d(1-z)^{\ell+2}}.
\]
This is precisely the bad-prefix argument from the endpoint theorem;
using the undivided polynomial includes rows containing zero entries.
Choose $\varepsilon=d^{-1/2}$. On the complementary good rows,
$\prod_{i\notin J}(1+zr_i)$ is squeezed between exponential moments
with parameters $z$ and $(1+\varepsilon)z$, as in that theorem.
Their common limit is $(1-z)^{-\ell-1}$.
The omitted mass tends to zero. Exponential tightness, uniform
integrability at a smaller positive parameter, and uniqueness of
moment generating functions identify the claimed Gamma limit.
All bounds depend only on $\ell,d,z,\varepsilon$, not on $J$.
\end{proof}

\begin{theorem}[Uniform profile]\label{profile}
For every fixed $L<\infty$,
\[
 \sup_{\substack{q:A(q)>0\\\lambda(q)\le L}}
 \max_{1\le j\le m}|m r_j(q)-\lambda(q)|\longrightarrow0.
\]
Equivalently, the same assertion holds with $r_j$ replaced by $1-p_j$.
It permits flat columns, unequal row weights, and different choices
of the representations at every $m$.
\end{theorem}
\begin{proof}
Choose any sequence of indices $j=j_m$. Put $\lambda_{-j}=\lambda-r_j$
on $A>0$. Under $\eta$, its law converges to $\operatorname{Exp}(1)$,
because $0\le r_j\le\max_i r_i\to0$ in probability. This convergence
is uniform in the choice of $j$.
By Lemma~\ref{gamma}, under $\nu_j$ the law of $\lambda_{-j}$ converges
to $\operatorname{Gamma}(2,1)$.

We claim that $r_j\to0$ in $\nu_j$ probability as well, assigning
$r_j=+\infty$ where $p_j=0$. Fix $\varepsilon>0$ and let
$D_j=\{r_j\ge\varepsilon\}$, an initial segment. Its $\eta$-mass is
at most $1/(m\varepsilon)$. For every fixed $B<\infty$, all rows of
$D_j$ that have positive $\nu_j$-mass must eventually satisfy
$\lambda_{-j}>B$. Otherwise take the last such row in $D_j$ with
$\lambda_{-j}\le B$. Since $\lambda_{-j}$ decreases with the row,
\[
 \eta(\lambda_{-j}>B)\le\eta(D_j)\longrightarrow0,
\]
contradicting its exponential limit. Rows where an odds ratio outside
$j$ is infinite have zero $\nu_j$-mass and can be omitted.
Tightness of the Gamma limit therefore gives
$\nu_j(D_j)\to0$. It follows that, after discarding zero-$A$ rows of
vanishing $\nu_j$-mass, the law of the full $\lambda$ under $\nu_j$
also converges to $\operatorname{Gamma}(2,1)$.
On $A>0$ the exact density relation is
\[
 d\nu_j=m r_j\,d\eta.
\]

Now choose any sequence of rows with $\lambda(q_m)\to x<\infty$.
For an interval $I$ lying strictly above $x$, monotonicity implies
$m r_j(q_m)\le\nu_j(\lambda\in I)/\eta(\lambda\in I)$ for all large
$m$. For an interval strictly below $x$, the inequality is reversed.
Every interval with positive length inside $(0,\infty)$ has positive
limiting $\eta$-mass, and the ratio tends to
\[
 \frac{\int_I t e^{-t}\,dt}{\int_I e^{-t}\,dt},
\]
which lies between the endpoints of $I$. Letting the bounding intervals
approach $x$ gives $m r_j(q_m)\to x$. For $x=0$, use only the upper
interval bound and nonnegativity.
If the displayed uniform assertion failed, a sequence of offending
indices and rows would have a subsequence with $\lambda(q_m)$ converging
in $[0,L]$, contradicting this conclusion.
Finally, uniformly on that region $m r_j$ is bounded and
\[
 0\le m r_j-m(1-p_j)=\frac{m r_j^2}{1+r_j}=O(1/m),
\]
proving the equivalent form.
\end{proof}

\begin{corollary}[The original tail mass is a linear coordinate]\label{tail}
For every fixed $L<\infty$, uniformly in $0\le x\le L$,
\[
 (m+1)\sum_{q:\,A(q)>0,\,\lambda(q)\le x}w_q\longrightarrow x.
\]
For any row $q$ with $A(q)>0$ and $\lambda(q)\le L$, put
$u_q=(m+1)\sum_{s=q}^K w_s$. Then, uniformly in these rows and columns,
\[
 \lambda(q)-u_q\longrightarrow0,\qquad
 m(1-p_j(q))-u_q\longrightarrow0.
\]
For equal weights, $u_q=(m+1)(K-q+1)/K$.
\end{corollary}
\begin{proof}
Theorem~\ref{profile} gives, uniformly for $\lambda\le L$,
\[
 -\log A=\sum_j\log(1+r_j)
 =\lambda+O\left((\max_jr_j)\lambda\right)=\lambda+o(1).
\]
Therefore the finite measures
\[
 (m+1)\sum_{q:\lambda(q)\le L}w_q\delta_{\lambda(q)}
 = A^{-1}\boldsymbol1_{[0,L]}\,d\eta
\]
converge weakly to Lebesgue measure on $[0,L]$, by the exponential
limit and the uniform approximation $A^{-1}=e^\lambda+o(1)$.
The cumulative functions converge uniformly because the limiting
measure has continuous cumulative function; extending the endpoint
slightly avoids any truncation ambiguity at $L$.
At a row of value $\lambda(q)$, its suffix mass lies between the
strict and non-strict cumulative masses at that value. Their difference
is a jump of the cumulative function, whose maximum tends to zero.
Thus $u_q-\lambda(q)\to0$ uniformly. The column conclusion follows
from Theorem~\ref{profile}.
\end{proof}

These are asymptotic structural statements. The independently proved
bound $K\ge c m\log m$ is stronger than the superlinear consequence
of the tail profile alone. This note supplies no further growth rate
for $K$.
\end{document}
